% =============================================================================
%  Chapter 58 — Building fall\_n on Windows (MSYS2/UCRT64 + GCC)
% =============================================================================

\chapter{Building \texttt{fall\_n} on Windows (MSYS2/UCRT64)}
\label{ch:windows-build}

This chapter documents the full procedure to configure, compile, and run
\texttt{fall\_n} on a 64-bit Windows machine using the MSYS2 UCRT64 toolchain.
Every step is reproducible from a clean Windows install and does not require
Visual Studio, WSL, or any other Unix compatibility layer beyond MSYS2 itself.

% =============================================================================
\section{Prerequisites and Motivations}
\label{sec:win-prereqs}

\texttt{fall\_n} targets C++23 and depends on several scientific libraries
(PETSc, Eigen, VTK, MS-MPI) that are either unavailable or extremely
inconvenient to build from source under the MSVC ABI.  The chosen toolchain is:

\begin{itemize}
  \item \textbf{MSYS2 UCRT64} --- a MinGW-w64 environment based on the
        Universal C Runtime (UCRT), which ships as a standard component of
        modern Windows.  All packages are installed via the \texttt{pacman}
        package manager, giving reproducible binary packages without local
        compilation.
  \item \textbf{GCC 15.2} --- full C++23 support including \texttt{<print>},
        \texttt{std::expected}, ranges, and concepts.
  \item \textbf{Ninja 1.13} + \textbf{CMake 4.2} --- fast incremental builds.
  \item \textbf{Microsoft MPI (MS-MPI)} --- the official MPI runtime for
        Windows, required by PETSc at link time and at runtime.
\end{itemize}

The companion library versions used and verified to work are summarised in
\autoref{tab:win-deps}.

\begin{table}[ht]
  \centering
  \caption{Dependency versions verified on Windows (MSYS2 UCRT64, March 2026).}
  \label{tab:win-deps}
  \begin{tabular}{llll}
    \toprule
    Library & MSYS2 package name & Version & Notes \\
    \midrule
    GCC           & \texttt{mingw-w64-ucrt-x86\_64-gcc}          & 15.2.0   & C++23 \\
    CMake         & \texttt{mingw-w64-ucrt-x86\_64-cmake}         & 4.2.3    & \\
    Ninja         & \texttt{mingw-w64-ucrt-x86\_64-ninja}         & 1.13.2   & \\
    pkg-config    & \texttt{mingw-w64-ucrt-x86\_64-pkgconf}       & 2.x      & \\
    Eigen         & \texttt{mingw-w64-ucrt-x86\_64-eigen3}        & 5.0.1    & header-only \\
    PETSc         & \texttt{mingw-w64-ucrt-x86\_64-petsc}         & 3.24.5   & double/MPI/opt \\
    VTK           & \texttt{mingw-w64-ucrt-x86\_64-vtk}           & 9.5.2    & \\
    MS-MPI SDK    & \texttt{mingw-w64-ucrt-x86\_64-msmpi}         & 10.1.1   & headers+stubs \\
    nlohmann-json & \texttt{mingw-w64-ucrt-x86\_64-nlohmann-json} & 3.12.0   & VTK dep \\
    fast\_float   & \texttt{mingw-w64-ucrt-x86\_64-fast\_float}   & 8.2.4    & VTK dep \\
    utf8cpp       & \texttt{mingw-w64-ucrt-x86\_64-utf8cpp}       & 4.0.6    & VTK dep \\
    MS-MPI runtime & \texttt{winget: Microsoft.msmpi}             & 10.1.12  & runtime DLL \\
    \bottomrule
  \end{tabular}
\end{table}

% =============================================================================
\section{Step 1 — Install MSYS2}
\label{sec:win-msys2}

Download the MSYS2 installer from \url{https://www.msys2.org} and run it.
The default installation path is \texttt{C:\textbackslash msys64}.
After installation, open the \textbf{MSYS2 UCRT64} terminal (not the plain
MSYS2 terminal) from the Start menu.

Update the base system before installing anything:
\begin{lstlisting}[language=bash, basicstyle=\ttfamily\small,
    keywordstyle=\bfseries\color{blue!70!black},
    commentstyle=\itshape\color{green!50!black},
    stringstyle=\color{red!60!black},
    frame=single, framerule=0.4pt, rulecolor=\color{gray!50},
    breaklines=true, showstringspaces=false]
pacman -Syu          # update database + core packages; close and reopen
pacman -Su           # finish updating remaining packages
\end{lstlisting}

% =============================================================================
\section{Step 2 — Install Build Tools and Dependencies}
\label{sec:win-pacman}

All subsequent commands are executed inside the \textbf{MSYS2 UCRT64} terminal.
Install the compiler toolchain and build tools first:

\begin{lstlisting}[language=bash, basicstyle=\ttfamily\small,
    keywordstyle=\bfseries\color{blue!70!black},
    commentstyle=\itshape\color{green!50!black},
    stringstyle=\color{red!60!black},
    frame=single, framerule=0.4pt, rulecolor=\color{gray!50},
    breaklines=true, showstringspaces=false]
pacman -S \
  mingw-w64-ucrt-x86_64-gcc \
  mingw-w64-ucrt-x86_64-cmake \
  mingw-w64-ucrt-x86_64-ninja \
  mingw-w64-ucrt-x86_64-pkgconf
\end{lstlisting}

Then install the scientific libraries and their transitive dependencies:

\begin{lstlisting}[language=bash, basicstyle=\ttfamily\small,
    keywordstyle=\bfseries\color{blue!70!black},
    commentstyle=\itshape\color{green!50!black},
    stringstyle=\color{red!60!black},
    frame=single, framerule=0.4pt, rulecolor=\color{gray!50},
    breaklines=true, showstringspaces=false]
pacman -S \
  mingw-w64-ucrt-x86_64-eigen3 \
  mingw-w64-ucrt-x86_64-petsc \
  mingw-w64-ucrt-x86_64-vtk \
  mingw-w64-ucrt-x86_64-msmpi \
  mingw-w64-ucrt-x86_64-nlohmann-json \
  mingw-w64-ucrt-x86_64-fast_float \
  mingw-w64-ucrt-x86_64-utf8cpp
\end{lstlisting}

\begin{remark}
The packages \texttt{nlohmann-json}, \texttt{fast\_float}, and \texttt{utf8cpp}
are direct dependencies of the MSYS2 VTK 9.5.2 package and must be installed
separately --- they are not pulled in automatically.  Omitting them causes
CMake's \texttt{find\_package(VTK)} to fail with missing component errors.
\end{remark}

\begin{remark}
The MSYS2 PETSc package uses the pkg-config name \texttt{petsc-dmo}
(double-precision, MPI-enabled, optimized build).  The legacy name
\texttt{PETSc} used in the original Linux \texttt{CMakeLists.txt} does not
exist; see \autoref{sec:win-cmake-changes} for the necessary
\texttt{pkg\_search\_module} adjustment.
\end{remark}

% =============================================================================
\section{Step 3 — Install the MS-MPI Runtime}
\label{sec:win-msmpi-runtime}

The \texttt{mingw-w64-ucrt-x86\_64-msmpi} pacman package provides the
MS-MPI \emph{SDK} (headers and import libraries) needed to compile and link
MPI programs.  However, the \emph{runtime DLL} (\texttt{msmpi.dll}) is
distributed separately by Microsoft and must be installed on every machine
that will \emph{run} executables linked against MS-MPI (including during
\texttt{ctest}).

Install the runtime using winget from a regular Windows PowerShell or Command
Prompt (not from inside MSYS2):

\begin{lstlisting}[language=bash, basicstyle=\ttfamily\small,
    keywordstyle=\bfseries\color{blue!70!black},
    commentstyle=\itshape\color{green!50!black},
    stringstyle=\color{red!60!black},
    frame=single, framerule=0.4pt, rulecolor=\color{gray!50},
    breaklines=true, showstringspaces=false]
winget install Microsoft.msmpi --accept-package-agreements
\end{lstlisting}

This installs \texttt{mpiexec.exe} to
\texttt{C:\textbackslash Program Files\textbackslash Microsoft MPI\textbackslash
Bin\textbackslash} and registers \texttt{msmpi.dll} in the Windows system
search path.  Without this step all PETSc-linked executables will abort at
startup with exit code \texttt{0xc0000135} (\textsc{dll not found}).

% =============================================================================
\section{Step 4 — Configure the CMake Build}
\label{sec:win-cmake-invoke}

Clone (or copy) the repository to a convenient location, e.g.\
\texttt{C:\textbackslash MyLibs\textbackslash fall\_n}.  All remaining
commands are executed inside the \textbf{MSYS2 UCRT64} terminal.

Create a build directory, configure, and generate:

\begin{lstlisting}[language=bash, basicstyle=\ttfamily\small,
    keywordstyle=\bfseries\color{blue!70!black},
    commentstyle=\itshape\color{green!50!black},
    stringstyle=\color{red!60!black},
    frame=single, framerule=0.4pt, rulecolor=\color{gray!50},
    breaklines=true, showstringspaces=false]
cd /c/MyLibs/fall_n
mkdir -p build && cd build
cmake .. --preset win-debug
\end{lstlisting}

The \texttt{win-debug} preset is defined in \texttt{CMakePresets.json} and
activates only on Windows hosts.  It selects the UCRT64 Ninja generator and
sets \texttt{CMAKE\_BUILD\_TYPE=Debug}.

A successful configuration ends with output similar to:
\begin{lstlisting}[basicstyle=\ttfamily\small, frame=single,
    framerule=0.4pt, rulecolor=\color{gray!50}, breaklines=true,
    showstringspaces=false]
-- Configuring done (3.6s)
-- Generating done (0.5s)
-- Build files have been written to: C:/MyLibs/fall_n/build
\end{lstlisting}

% =============================================================================
\section{Step 5 — Compile}
\label{sec:win-build}

From the build directory, invoke Ninja:

\begin{lstlisting}[language=bash, basicstyle=\ttfamily\small,
    keywordstyle=\bfseries\color{blue!70!black},
    commentstyle=\itshape\color{green!50!black},
    stringstyle=\color{red!60!black},
    frame=single, framerule=0.4pt, rulecolor=\color{gray!50},
    breaklines=true, showstringspaces=false]
ninja -j8         # build everything (main + seismic + all 48 tests)
ninja -j8 fall_n  # build only the main executable
\end{lstlisting}

On a modern 8-core machine the full build (including the Matplot++ dependency
fetched via \texttt{FetchContent}) takes approximately 4--6 minutes.  The
following executables are produced in the build directory:

\begin{itemize}
  \item \texttt{fall\_n.exe} --- main structural analysis executable;
  \item \texttt{fall\_n\_seismic.exe} --- seismic analysis entry point;
  \item 48 test executables (\texttt{fall\_n\_*\_test.exe}).
\end{itemize}

% =============================================================================
\section{Step 6 — Run the Tests}
\label{sec:win-tests}

\begin{lstlisting}[language=bash, basicstyle=\ttfamily\small,
    keywordstyle=\bfseries\color{blue!70!black},
    commentstyle=\itshape\color{green!50!black},
    stringstyle=\color{red!60!black},
    frame=single, framerule=0.4pt, rulecolor=\color{gray!50},
    breaklines=true, showstringspaces=false]
ctest -L unit -j4            # fast unit tests only (~4 s)
ctest -L 'unit|integration' -j4
ctest -j4                    # all 48 tests
\end{lstlisting}

Expected results on a clean install: \textbf{44/48 tests pass}.  The 4
remaining failures are not related to the Windows port; they are caused by
missing mesh and seismic data files that are not included in the repository:

\begin{itemize}
  \item \texttt{seismic\_infrastructure}: requires
        \texttt{data/input/earthquakes/el\_centro\_1940\_ns.dat};
  \item \texttt{nonlinear\_validation}, \texttt{vertical\_slice\_gate}:
        require \texttt{tests/validation\_cube.msh};
  \item \texttt{boundary\_conditions}: requires
        \texttt{data/input/Beam\_LagCell\_Ord1\_30x6x3.msh}.
\end{itemize}

% =============================================================================
\section{CMake Modifications for Windows Compatibility}
\label{sec:win-cmake-changes}

This section documents every change that was made to \texttt{CMakeLists.txt}
and \texttt{CMakePresets.json} relative to the original Linux build system.
All changes are guarded by platform conditions so the file remains fully
functional on Linux without modification.

% -----------------------------------------------------------------------------
\subsection{MinGW-specific compiler and linker flags}

GCC on MinGW-w64 requires two flags that are not needed on Linux:

\begin{enumerate}
  \item \textbf{\texttt{-lstdc++exp}} --- The C++23 \texttt{<print>} header
        (\texttt{std::print}, \texttt{std::println}) is implemented in a
        separate experimental static library on MinGW.  Without this flag all
        translation units that include \texttt{<print>} fail at link time with
        undefined references to \texttt{\_ZSt5printB5cxx11\ldots}.

  \item \textbf{\texttt{-Wa,-mbig-obj}} --- The PE/COFF object file format
        (used on Windows) limits the number of sections in a single
        \texttt{.obj} file to 65\,535.  Template-heavy translation units in
        \texttt{fall\_n} (particularly those including PETSc and VTK headers)
        exceed this limit and cause the assembler to abort with the error
        \emph{``too many sections (73659 > 65535)''}.  This flag enables the
        \emph{big-obj} format which raises the limit to over four billion
        sections.
\end{enumerate}

The corresponding CMake block added to \texttt{CMakeLists.txt} is:
\begin{lstlisting}[language=C++, basicstyle=\ttfamily\small,
    keywordstyle=\bfseries\color{blue!70!black},
    commentstyle=\itshape\color{green!50!black},
    stringstyle=\color{red!60!black},
    frame=single, framerule=0.4pt, rulecolor=\color{gray!50},
    breaklines=true, showstringspaces=false]
if(MINGW AND CMAKE_CXX_COMPILER_ID STREQUAL "GNU")
  link_libraries(stdc++exp)
  add_compile_options(-Wa,-mbig-obj)

  # VTK DLLs on MinGW do not export inline template member functions.
  # Predefine these guards to skip the extern template declarations
  # and allow the compiler to instantiate the functions locally.
  add_compile_definitions(
    VTK_AOS_DATA_ARRAY_TEMPLATE_EXTERN
    VTK_GDA_TEMPLATE_EXTERN
  )
endif()
\end{lstlisting}

% -----------------------------------------------------------------------------
\subsection{VTK extern-template linking fix}
\label{ssec:win-vtk-extern}

This is the most subtle Windows-specific issue.  VTK uses \emph{explicit
template instantiation} to reduce compile times: the class template
\texttt{vtkAOSDataArrayTemplate<T>} (and its ancestor
\texttt{vtkGenericDataArray<>}) is fully instantiated inside the
\texttt{libvtkCommonCore.dll} shared library, while all consumer translation
units use \texttt{extern template} declarations (gated by the macro
\texttt{VTK\_USE\_EXTERN\_TEMPLATE}, unconditionally defined in
\texttt{vtkCompiler.h}) to suppress local re-instantiation.

On Linux and macOS this mechanism works perfectly: the shared library exports
the symbols and the dynamic linker resolves them at load time.  On Windows
with MinGW, however, the \texttt{extern template} suppression interacts poorly
with the PE/COFF \texttt{\_\_declspec(dllimport)} mechanism.  In particular,
\emph{inline methods} such as

\begin{lstlisting}[language=C++, basicstyle=\ttfamily\small,
    frame=single, framerule=0.4pt, rulecolor=\color{gray!50},
    breaklines=true, showstringspaces=false]
// vtkAOSDataArrayTemplate.h (line 87)
void SetValue(vtkIdType valueIdx, ValueType value) {
  this->Buffer->GetBuffer()[valueIdx] = value;
}
\end{lstlisting}

are \emph{already inlined} into the DLL's callers at the time the DLL is
built; they are \emph{not emitted as distinct exported symbols}.  Yet
\texttt{extern template class VTKCOMMONCORE\_EXPORT
vtkAOSDataArrayTemplate<T>;} instructs the consumer translation unit to
\emph{not} instantiate the template locally and to expect the symbol from the
DLL import library --- a symbol that does not exist.  The linker then fails
with:
\begin{lstlisting}[basicstyle=\ttfamily\small, frame=single,
    framerule=0.4pt, rulecolor=\color{gray!50}, breaklines=true,
    showstringspaces=false]
undefined reference to
  `vtkAOSDataArrayTemplate<long long>::SetValue(long long, long long)'
undefined reference to
  `vtkAOSDataArrayTemplate<int>::SetValue(int, int)'
undefined reference to
  `vtkAOSDataArrayTemplate<double>::SetValue(double, double)'
\end{lstlisting}

The fix is to predefine the guard macros \texttt{VTK\_AOS\_DATA\_ARRAY\_TEMPLATE\_EXTERN}
and \texttt{VTK\_GDA\_TEMPLATE\_EXTERN} before any VTK header is processed.
These macros wrap the \texttt{extern template} declarations inside
\texttt{\#ifndef} guards:
\begin{lstlisting}[language=C++, basicstyle=\ttfamily\small,
    frame=single, framerule=0.4pt, rulecolor=\color{gray!50},
    breaklines=true, showstringspaces=false]
// vtkAOSDataArrayTemplate.h (abbreviated)
#ifdef VTK_USE_EXTERN_TEMPLATE
#ifndef VTK_AOS_DATA_ARRAY_TEMPLATE_EXTERN   // <-- guard
#define VTK_AOS_DATA_ARRAY_TEMPLATE_EXTERN
extern template class VTKCOMMONCORE_EXPORT vtkAOSDataArrayTemplate<double>;
// ... 12 more explicit instantiations ...
#endif
#endif
\end{lstlisting}

By predeclaring \texttt{VTK\_AOS\_DATA\_ARRAY\_TEMPLATE\_EXTERN} (and the
analogous \texttt{VTK\_GDA\_TEMPLATE\_EXTERN} for the base class), the
\texttt{extern template} block is skipped and the compiler instantiates the
inline methods locally from the \texttt{.h} + \texttt{.txx} pair shipped with
the package.  This is the correct semantics on MinGW: small inline methods
should never have been placed behind DLL-import stubs in the first place.

% -----------------------------------------------------------------------------
\subsection{PETSc C++ macro fallbacks}

The MSYS2 PETSc 3.24.5 package was compiled as a pure-C library
(\texttt{PETSC\_CLANGUAGE\_C~=~1}).  As a result,
two macros required by the C++ PETSc headers are absent from
\texttt{petscconf.h}:

\begin{itemize}
  \item \texttt{PETSC\_FUNCTION\_NAME\_CXX} --- the C++ equivalent of
        \texttt{\_\_FUNCTION\_\_}, used in error-handling macros such as
        \texttt{PetscCall} for stack-trace annotation.
  \item \texttt{PETSC\_CXX\_RESTRICT} --- the C++ spelling of the
        \texttt{restrict} qualifier.
\end{itemize}

Missing either macro causes compile errors of the form \emph{``PETSC\_FUNCTION\_NAME\_CXX was not declared''}.  The fix uses CMake's
\texttt{CheckSymbolExists} to detect their absence and defines fallbacks:

\begin{lstlisting}[language=C++, basicstyle=\ttfamily\small,
    keywordstyle=\bfseries\color{blue!70!black},
    commentstyle=\itshape\color{green!50!black},
    stringstyle=\color{red!60!black},
    frame=single, framerule=0.4pt, rulecolor=\color{gray!50},
    breaklines=true, showstringspaces=false]
include(CheckSymbolExists)
set(CMAKE_REQUIRED_INCLUDES ${PETSC_INCLUDE_DIRS})
check_symbol_exists(PETSC_FUNCTION_NAME_CXX "petscconf.h"
                    _HAVE_PETSC_FN_CXX)
if(NOT _HAVE_PETSC_FN_CXX)
  add_compile_definitions(
    PETSC_FUNCTION_NAME_CXX=__PRETTY_FUNCTION__)
endif()
check_symbol_exists(PETSC_CXX_RESTRICT "petscconf.h"
                    _HAVE_PETSC_CXX_RESTRICT)
if(NOT _HAVE_PETSC_CXX_RESTRICT)
  add_compile_definitions(PETSC_CXX_RESTRICT=__restrict__)
endif()
\end{lstlisting}

% -----------------------------------------------------------------------------
\subsection{PETSc pkg-config name}

The original \texttt{CMakeLists.txt} searched for PETSc via
\texttt{pkg\_search\_module(PETSC REQUIRED IMPORTED\_TARGET PETSc)}.  The
MSYS2 package registers the library under the name \texttt{petsc-dmo} (double,
MPI, optimized).  The search was widened to try multiple names:

\begin{lstlisting}[language=C++, basicstyle=\ttfamily\small,
    frame=single, framerule=0.4pt, rulecolor=\color{gray!50},
    breaklines=true, showstringspaces=false]
pkg_search_module(PETSC REQUIRED IMPORTED_TARGET
                  PETSc petsc petsc-dmo petsc-dso)
\end{lstlisting}

% -----------------------------------------------------------------------------
\subsection{Source directory macro}

Hardcoded absolute Unix paths of the form
\texttt{/home/sechavarriam/MyLibs/fall\_n/} existed in \texttt{main.cpp},
\texttt{main\_seismic.cpp}, and thirteen test files.  These were replaced with
a compile-time macro that resolves to the actual source tree regardless of
platform or user name:

\begin{lstlisting}[language=C++, basicstyle=\ttfamily\small,
    frame=single, framerule=0.4pt, rulecolor=\color{gray!50},
    breaklines=true, showstringspaces=false]
# CMakeLists.txt
add_compile_definitions(
  FALL_N_SOURCE_DIR="${CMAKE_SOURCE_DIR}")
\end{lstlisting}

Usage in source code:
\begin{lstlisting}[language=C++, basicstyle=\ttfamily\small,
    frame=single, framerule=0.4pt, rulecolor=\color{gray!50},
    breaklines=true, showstringspaces=false]
#ifdef FALL_N_SOURCE_DIR
  static const std::string BASE =
      std::string(FALL_N_SOURCE_DIR) + "/";
#else
  static const std::string BASE = "./";
#endif
\end{lstlisting}

% -----------------------------------------------------------------------------
\subsection{Windows CMake preset}

A \texttt{win-debug} preset was added to \texttt{CMakePresets.json}:
\begin{lstlisting}[language=C++, basicstyle=\ttfamily\small,
    frame=single, framerule=0.4pt, rulecolor=\color{gray!50},
    breaklines=true, showstringspaces=false]
{
  "name": "win-debug",
  "displayName": "Windows Debug (MSYS2 UCRT64)",
  "generator": "Ninja",
  "binaryDir": "${sourceDir}/build",
  "cacheVariables": { "CMAKE_BUILD_TYPE": "Debug" },
  "condition": {
    "type": "equals",
    "lhs": "${hostSystemName}",
    "rhs": "Windows"
  }
}
\end{lstlisting}

% =============================================================================
\section{Source-Code Portability Fixes}
\label{sec:win-src-fixes}

Besides the build-system changes, several source files required modifications
to compile cleanly with GCC on Windows.

% -----------------------------------------------------------------------------
\subsection{\texttt{ushort} is not a standard type}

The file \texttt{src/utils/index.hh} used \texttt{ushort} as a template
non-type parameter:
\begin{lstlisting}[language=C++, basicstyle=\ttfamily\small,
    frame=single, framerule=0.4pt, rulecolor=\color{gray!50},
    breaklines=true, showstringspaces=false]
template <ushort... N>  // non-standard: ushort is a POSIX extension
\end{lstlisting}
GCC on MinGW does not declare \texttt{ushort} in the global namespace.
The fix replaces it with the standard \texttt{unsigned short}:
\begin{lstlisting}[language=C++, basicstyle=\ttfamily\small,
    frame=single, framerule=0.4pt, rulecolor=\color{gray!50},
    breaklines=true, showstringspaces=false]
template <unsigned short... N>
\end{lstlisting}

% -----------------------------------------------------------------------------
\subsection{\texttt{\%zu} format specifier in \texttt{PetscPrintf}}

The MSYS2 PETSc runtime's \texttt{PetscPrintf} does not recognise the
\texttt{\%zu} format specifier for \texttt{std::size\_t} on this platform.
Occurrences were replaced with \texttt{\%d} and an explicit \texttt{(int)} cast:
\begin{lstlisting}[language=C++, basicstyle=\ttfamily\small,
    frame=single, framerule=0.4pt, rulecolor=\color{gray!50},
    breaklines=true, showstringspaces=false]
// Before
PetscPrintf(PETSC_COMM_WORLD,
    "T_measured = %.6f (from %zu crossings)\n",
    T_measured, crossings.size());

// After
PetscPrintf(PETSC_COMM_WORLD,
    "T_measured = %.6f (from %d crossings)\n",
    T_measured, (int)crossings.size());
\end{lstlisting}

% -----------------------------------------------------------------------------
\subsection{Unused-variable warning in tests}

GCC 15 with \texttt{-Werror=unused-but-set-variable} rejected a variable in
\texttt{tests/test\_kinematic\_policy.cpp} that was assigned but never read
after the assignment.  The fix adds \texttt{[[maybe\_unused]]}:
\begin{lstlisting}[language=C++, basicstyle=\ttfamily\small,
    frame=single, framerule=0.4pt, rulecolor=\color{gray!50},
    breaklines=true, showstringspaces=false]
[[maybe_unused]] Eigen::Vector<double, 3> E_from_B = gp.B * u_e;
\end{lstlisting}

% =============================================================================
\section{Runtime PATH Considerations}
\label{sec:win-runtime-path}

The UCRT64 runtime libraries (\texttt{libvtkCommonCore.dll},
\texttt{libpetsc-dmo.dll}, \texttt{libgcc\_s\_seh-1.dll}, etc.) reside in
\texttt{C:\textbackslash msys64\textbackslash ucrt64\textbackslash bin}.  This
directory is automatically on \texttt{PATH} inside MSYS2 terminal sessions.

When executables are launched \emph{outside} MSYS2 (e.g.\ from PowerShell,
Explorer, or a CI runner) they will fail with \texttt{0xc0000135} unless the
UCRT64 bin path is added to the system or user \texttt{PATH}, or the DLLs are
copied alongside the executables.

For CTest invocations the recommended approach is to run from within the MSYS2
UCRT64 terminal, where the path is already configured:
\begin{lstlisting}[language=bash, basicstyle=\ttfamily\small,
    keywordstyle=\bfseries\color{blue!70!black},
    commentstyle=\itshape\color{green!50!black},
    stringstyle=\color{red!60!black},
    frame=single, framerule=0.4pt, rulecolor=\color{gray!50},
    breaklines=true, showstringspaces=false]
# Open MSYS2 UCRT64 terminal, then:
cd /c/MyLibs/fall_n/build
ctest -L unit -j4 --output-on-failure
\end{lstlisting}

Alternatively, add the UCRT64 bin directory permanently to the Windows system
environment variables via \textit{System Properties $\to$ Environment Variables
$\to$ System Variables $\to$ Path}.

% =============================================================================
\section{Summary of All Changes}
\label{sec:win-summary}

\autoref{tab:win-changes} lists every file that was modified and the
corresponding reason.

\begin{table}[ht]
  \centering
  \caption{Files modified for Windows/MinGW compatibility.}
  \label{tab:win-changes}
  \begin{tabular}{p{6.2cm} p{8cm}}
    \toprule
    File & Change \\
    \midrule
    \texttt{CMakeLists.txt}
      & MinGW flags (\texttt{-lstdc++exp},
        \texttt{-Wa,-mbig-obj}); PETSc fallback macros;
        \texttt{pkg\_search\_module} widened; VTK extern-template guards;
        \texttt{FALL\_N\_SOURCE\_DIR} macro; compiler flag guard for MSVC
        vs.\ GCC. \\[3pt]
    \texttt{CMakePresets.json}
      & Added \texttt{win-debug} preset. \\[3pt]
    \texttt{main.cpp}
      & Replaced hardcoded \texttt{/home/\ldots} path with
        \texttt{FALL\_N\_SOURCE\_DIR} macro. \\[3pt]
    \texttt{main\_seismic.cpp}
      & Same as \texttt{main.cpp}. \\[3pt]
    \texttt{src/utils/index.hh}
      & \texttt{ushort} $\to$ \texttt{unsigned short}. \\[3pt]
    13 test files under \texttt{tests/}
      & Replaced hardcoded Unix paths with \texttt{FALL\_N\_SOURCE\_DIR};
        replaced \texttt{/tmp/} with
        \texttt{std::filesystem::temp\_directory\_path()}. \\[3pt]
    \texttt{tests/test\_kinematic\_policy.cpp}
      & Added \texttt{[[maybe\_unused]]} to suppress GCC 15
        unused-variable warning. \\[3pt]
    \texttt{tests/test\_phase5\_dynamic\_verification.cpp}
      & Replaced \texttt{\%zu} with \texttt{\%d} in
        \texttt{PetscPrintf} format string. \\
    \bottomrule
  \end{tabular}
\end{table}
